\documentclass{article}
\usepackage{amsmath,amssymb,amsthm}
\usepackage{algorithm}
\usepackage{algpseudocode}
\usepackage{hyperref}
\usepackage{enumitem}
\usepackage{graphicx}
\usepackage{bm}

\newtheorem{theorem}{Theorem}
\newtheorem{lemma}{Lemma}
\newtheorem{assumption}{Assumption}

\begin{document}

\section*{Cyclic Stochastic Gradient Descent (Cyclic SGD)}

\section*{Mathematical Formulation}
Let $f:\mathbb{R}^{d}\to\mathbb{R}$ be the (possibly non‑convex) objective function and let
$\{z_{t}\}_{t\ge 0}$ be i.i.d.\ data samples.  
Denote by 
\[
g_{t}\;:=\;\nabla f(w_{t};z_{t})
\]
a stochastic gradient evaluated at the current iterate $w_{t}\in\mathbb{R}^{d}$.

The learning‑rate schedule is \emph{cyclic} with period $M\in\mathbb{N}$, lower bound
$\eta_{\min}>0$ and upper bound $\eta_{\max}> \eta_{\min}$.  
For $t\ge 0$ define the intra‑cycle index 
\[
\kappa_{t}\;:=\;t\bmod (2M)\in\{0,\dots ,2M-1\}.
\]
The \textbf{triangular} cyclic step size is
\[
\boxed{\;
\eta_{t}\;=\;\eta_{\min}
      \;+\;(\eta_{\max}-\eta_{\min})\;
        \Bigl(1-\Bigl|\frac{\kappa_{t}}{M}-1\Bigr|\Bigr)
\;}
\tag{1}
\]
which linearly increases from $\eta_{\min}$ to $\eta_{\max}$ in the first $M$ steps
and then linearly decreases back to $\eta_{\min}$ in the next $M$ steps.

The update rule of Cyclic SGD is
\[
\boxed{\; w_{t+1}\;=\;w_{t}\;-\;\eta_{t}\,g_{t}\; } \qquad t=0,1,2,\dots
\tag{2}
\]

\section*{Pseudocode}
\begin{algorithm}[H]
\caption{Cyclic Stochastic Gradient Descent}
\begin{algorithmic}[1]
\State \textbf{Input:} initial point $w_{0}\in\mathbb{R}^{d}$,
period $M$, bounds $\eta_{\min},\eta_{\max}$,
maximal iteration $T$.
\For{$t=0$ \textbf{to} $T-1$}
    \State Compute intra‑cycle index $\kappa_{t}=t\bmod (2M)$.
    \State Set learning rate 
    \[
    \eta_{t}\gets \eta_{\min}+(\eta_{\max}-\eta_{\min})\Bigl(1-\Bigl|\frac{\kappa_{t}}{M}-1\Bigr|\Bigr).
    \] \Comment{triangular schedule (1)}
    \State Sample a data point $z_{t}$ and compute stochastic gradient 
    $g_{t}\gets\nabla f(w_{t};z_{t})$.
    \State Update $w_{t+1}\gets w_{t}-\eta_{t}\,g_{t}$. \Comment{update (2)}
\EndFor
\State \textbf{Return} $w_{T}$ (or the averaged iterate $\bar w_{T}=\frac1T\sum_{t=0}^{T-1}w_{t}$).
\end{algorithmic}
\end{algorithm}

\section*{Dry Run Example}
Consider the simple quadratic $f(w)=(w-3)^{2}$, $w\in\mathbb{R}$.  
We take $w_{0}=0$, period $M=2$, $\eta_{\min}=0.05$, $\eta_{\max}=0.15$ and
assume exact gradients (i.e., $g_{t}=\nabla f(w_{t})$).  
The triangular schedule (1) yields:

\begin{center}
\begin{tabular}{c|c|c|c}
$t$ & $\kappa_{t}$ & $\eta_{t}$ & $w_{t}$ \\ \hline
0 & $0$ & $0.05$ & $0$ \\
1 & $1$ & $0.10$ & $0$ \\
2 & $2$ & $0.15$ & $?$ \\
3 & $3$ & $0.10$ & $?$ \\
\end{tabular}
\end{center}

\begin{enumerate}[label={\bf Step \arabic*:}, leftmargin=*, nosep]
\item $t=0$: $\nabla f(w_{0})=2(w_{0}-3)=-6$.
   \[
   w_{1}=w_{0}-\eta_{0}\nabla f(w_{0})=0-0.05(-6)=0.30.
   \]
   Objective: $f(w_{1})=(0.30-3)^{2}=7.29$.
\item $t=1$: $\nabla f(w_{1})=2(0.30-3)=-5.40$.
   \[
   w_{2}=w_{1}-\eta_{1}\nabla f(w_{1})=0.30-0.10(-5.40)=0.84.
   \]
   Objective: $f(w_{2})=(0.84-3)^{2}=4.66$.
\item $t=2$: $\nabla f(w_{2})=2(0.84-3)=-4.32$.
   \[
   w_{3}=w_{2}-\eta_{2}\nabla f(w_{2})=0.84-0.15(-4.32)=1.49.
   \]
   Objective: $f(w_{3})=(1.49-3)^{2}=2.28$.
\end{enumerate}
The example illustrates how the step size cycles and the iterates move
towards the minimiser $w^{\star}=3$.

\newpage
\section*{Assumptions}
\begin{assumption}[Smoothness]\label{ass:smooth}
$f$ is $L$‑smooth, i.e.
\[
\|\nabla f(x)-\nabla f(y)\|\le L\|x-y\|\quad\forall x,y\in\mathbb{R}^{d}.
\tag{A1}
\]
Equivalently,
\[
f(y)\le f(x)+\langle\nabla f(x),y-x\rangle+\frac{L}{2}\|y-x\|^{2}.
\tag{A1'}
\]
\end{assumption}

\begin{assumption}[Unbiased Stochastic Gradient]\label{ass:unbiased}
For every $t$,
\[
\mathbb{E}[\,g_{t}\mid w_{t}\,]=\nabla f(w_{t}). 
\tag{A2}
\]
\end{assumption}

\begin{assumption}[Bounded Variance]\label{ass:variance}
There exists $\sigma^{2}\ge0$ such that
\[
\mathbb{E}\bigl[\|g_{t}-\nabla f(w_{t})\|^{2}\mid w_{t}\bigr]\le\sigma^{2},
\qquad\forall t.
\tag{A3}
\]
\end{assumption}

\begin{assumption}[Learning‑Rate Bounds]\label{ass:eta}
The cyclic schedule satisfies
\[
0<\eta_{\min}\le\eta_{t}\le\eta_{\max}\qquad\forall t,
\]
and the period $M$ is a fixed positive integer.
\tag{A4}
\end{assumption}

\section*{Main Convergence Theorem}
\begin{theorem}[Convergence of Cyclic SGD]\label{thm:main}
Under Assumptions \ref{ass:smooth}–\ref{ass:eta} and for any
$T\ge 1$, let $\bar w_{T}:=\frac{1}{T}\sum_{t=0}^{T-1}w_{t}$.
Define the average learning rate
\[
\bar\eta_{T}:=\frac{1}{T}\sum_{t=0}^{T-1}\eta_{t}.
\]
Then the following bound on the expected squared gradient norm holds:
\[
\frac{1}{T}\sum_{t=0}^{T-1}\mathbb{E}\!\bigl[\|\nabla f(w_{t})\|^{2}\bigr]
\;\le\;
\frac{2\bigl(f(w_{0})-f^{\star}\bigr)}{L\,\bar\eta_{T}T}
\;+\;
\frac{L\,\eta_{\max}\sigma^{2}}{2}.
\tag{3}
\]
Consequently, if the maximal step size satisfies
$\eta_{\max}=O(T^{-1/2})$, then
\[
\frac{1}{T}\sum_{t=0}^{T-1}\mathbb{E}\!\bigl[\|\nabla f(w_{t})\|^{2}\bigr]
=O\!\bigl(T^{-1/2}\bigr).
\]
\end{theorem}

\section*{Theoretical Properties}
\begin{enumerate}[label=\arabic*. , leftmargin=*]
\item \textbf{Stability.} Because every $\eta_{t}$ stays within $[\eta_{\min},\eta_{\max}]$
the iterates remain bounded whenever the level set $\{w\mid f(w)\le f(w_{0})\}$ is bounded,
a standard consequence of $L$‑smoothness.

\item \textbf{Hyper‑parameter Sensitivity.} The bound (3) depends on
$\eta_{\max}$ linearly in the variance term and on the harmonic mean
$\bar\eta_{T}$ in the deterministic term. Increasing $\eta_{\max}$ speeds up
the descent (larger $\bar\eta_{T}$) but also inflates the variance penalty.
The cyclic pattern does not worsen the worst‑case bound compared to a
constant step size $\eta=\bar\eta_{T}$.

\item \textbf{Comparison to Baselines.}
\begin{itemize}
\item \emph{Constant SGD:} Setting $\eta_{t}\equiv\eta$ recovers the classical bound
$\frac{2(f(w_{0})-f^{\star})}{L\eta T}+ \frac{L\eta\sigma^{2}}{2}$.
\item \emph{Cyclic SGD:} The bound (3) is identical after replacing $\eta$ by the
average $\bar\eta_{T}$ and the variance term by $L\eta_{\max}\sigma^{2}/2$.
Thus cyclic schedules achieve the same $O(T^{-1/2})$ rate while empirically
often explore the landscape more aggressively.
\end{itemize}
\end{enumerate}

\section*{Discussion}
The theorem shows that the cyclic modulation of the learning rate does not
damage the worst‑case convergence guarantees for smooth non‑convex objectives.
The analysis relies only on the fact that the step sizes are uniformly bounded;
their deterministic periodic pattern does not interfere with the telescoping
argument.  In practice, the temporary large step sizes ($\eta_{\max}$) help the
algorithm escape shallow basins, while the subsequent reduction to $\eta_{\min}$
provides a stabilising effect, a phenomenon that is consistent with the
theoretical decomposition into a deterministic descent term and a stochastic
variance term.

\newpage
\section*{Preliminaries (Definitions and Lemmas)}
\begin{lemma}[Descent Lemma]\label{lem:descent}
If $f$ is $L$‑smooth then for any $x,y\in\mathbb{R}^{d}$,
\[
f(y)\le f(x)+\langle\nabla f(x),y-x\rangle+\frac{L}{2}\|y-x\|^{2}.
\]
\end{lemma}
\begin{proof}
Directly from Assumption \ref{ass:smooth} (inequality (A1')). \hfill\qedsymbol
\end{proof}

\begin{lemma}[Variance Decomposition]\label{lem:var}
For any random vector $g$ with $\mathbb{E}[g]=\mu$,
\[
\mathbb{E}\bigl[\|g\|^{2}\bigr]=\|\mu\|^{2}+\mathbb{E}\bigl[\|g-\mu\|^{2}\bigr].
\]
\end{lemma}
\begin{proof}
Expand $\|g\|^{2}=\|g-\mu+\mu\|^{2}=\|g-\mu\|^{2}+2\langle g-\mu,\mu\rangle+\|\mu\|^{2}$,
take expectation, the cross term vanishes. \hfill\qedsymbol
\end{proof}

\section*{Main Proof (Step‑by‑Step Derivation)}
We prove Theorem \ref{thm:main}.  Throughout the proof $\mathbb{E}_{t}[\cdot]$
denotes $\mathbb{E}[\cdot\mid w_{t}]$.

\begin{enumerate}[label=[Step \arabic*], leftmargin=*]
\item \textbf{Apply the Descent Lemma.}  
Using Lemma \ref{lem:descent} with $x=w_{t}$, $y=w_{t+1}=w_{t}-\eta_{t}g_{t}$,
\[
f(w_{t+1})\le f(w_{t})-\eta_{t}\langle\nabla f(w_{t}),g_{t}\rangle
+\frac{L}{2}\eta_{t}^{2}\|g_{t}\|^{2}.
\tag{4}
\]

\item \textbf{Take conditional expectation.}  
Condition on $w_{t}$ and use unbiasedness (A2):
\[
\mathbb{E}_{t}[\,\langle\nabla f(w_{t}),g_{t}\rangle\,]
= \|\nabla f(w_{t})\|^{2}.
\tag{5}
\]

\item \textbf{Bound the second moment of $g_{t}$.}  
From Lemma \ref{lem:var} and the variance bound (A3):
\[
\mathbb{E}_{t}[\|g_{t}\|^{2}]
= \|\nabla f(w_{t})\|^{2}+\mathbb{E}_{t}[\|g_{t}-\nabla f(w_{t})\|^{2}]
\le \|\nabla f(w_{t})\|^{2}+\sigma^{2}.
\tag{6}
\]

\item \textbf{Insert (5)–(6) into (4).}  
Taking $\mathbb{E}_{t}$ on both sides of (4) yields
\[
\mathbb{E}_{t}[f(w_{t+1})]\le
f(w_{t})-\eta_{t}\|\nabla f(w_{t})\|^{2}
+\frac{L}{2}\eta_{t}^{2}\bigl(\|\nabla f(w_{t})\|^{2}+\sigma^{2}\bigr).
\tag{7}
\]

\item \textbf{Rearrange to isolate the gradient term.}  
Move the deterministic gradient term to the left:
\[
\bigl(\eta_{t}-\tfrac{L}{2}\eta_{t}^{2}\bigr)\|\nabla f(w_{t})\|^{2}
\le f(w_{t})-\mathbb{E}_{t}[f(w_{t+1})]
+\tfrac{L}{2}\eta_{t}^{2}\sigma^{2}.
\tag{8}
\]

\item \textbf{Use the bound $\eta_{t}\le \eta_{\max}$ to simplify the coefficient.}  
Since $0<\eta_{t}\le\eta_{\max}$ and $L\eta_{\max}\le 1$ (a standard step‑size
restriction; we can always enforce it by scaling $\eta_{\max}$), we have
\[
\eta_{t}-\tfrac{L}{2}\eta_{t}^{2}\;\ge\;\tfrac{1}{2}\eta_{t}.
\tag{9}
\]
(Indeed, $\eta_{t}(1-\frac{L}{2}\eta_{t})\ge\frac12\eta_{t}$ iff
$1-\frac{L}{2}\eta_{t}\ge\frac12$, i.e. $L\eta_{t}\le1$.)

\item \textbf{Combine (8) and (9).}
\[
\frac{1}{2}\eta_{t}\,\|\nabla f(w_{t})\|^{2}
\le f(w_{t})-\mathbb{E}_{t}[f(w_{t+1})]
+\frac{L}{2}\eta_{t}^{2}\sigma^{2}.
\tag{10}
\]

\item \textbf{Take full expectation and sum over $t=0,\dots,T-1$.}  
Define $\mathbb{E}$ as expectation over all randomness. Summing (10) yields
\[
\frac{1}{2}\sum_{t=0}^{T-1}\eta_{t}\,
\mathbb{E}\!\bigl[\|\nabla f(w_{t})\|^{2}\bigr]
\le \mathbb{E}[f(w_{0})]-\mathbb{E}[f(w_{T})]
+\frac{L\sigma^{2}}{2}\sum_{t=0}^{T-1}\eta_{t}^{2}.
\tag{11}
\]

\item \textbf{Lower bound the left‑hand side by the average learning rate.}  
Since $\eta_{t}\ge\eta_{\min}>0$,
\[
\frac{1}{2}\sum_{t=0}^{T-1}\eta_{t}\,
\mathbb{E}\!\bigl[\|\nabla f(w_{t})\|^{2}\bigr]
\ge
\frac{\eta_{\min}}{2}\sum_{t=0}^{T-1}
\mathbb{E}\!\bigl[\|\nabla f(w_{t})\|^{2}\bigr].
\tag{12}
\]
However, to obtain a bound expressed via the \emph{average} step size,
divide both sides of (11) by $T$ and use $\bar\eta_{T}=\frac1T\sum_{t=0}^{T-1}\eta_{t}$:
\[
\frac{1}{T}\sum_{t=0}^{T-1}\mathbb{E}\!\bigl[\|\nabla f(w_{t})\|^{2}\bigr]
\le
\frac{2\bigl(\mathbb{E}[f(w_{0})]-\mathbb{E}[f(w_{T})]\bigr)}{L\,\bar\eta_{T}T}
+\frac{L\,\eta_{\max}\sigma^{2}}{2}.
\tag{13}
\]
The inequality $\sum_{t}\eta_{t}^{2}\le \eta_{\max}\sum_{t}\eta_{t}
=\eta_{\max}T\bar\eta_{T}$ was used in the last term.

\item \textbf{Replace $\mathbb{E}[f(w_{T})]$ by the optimal value $f^{\star}$.}  
Since $f^{\star}\le f(w)$ for all $w$, we obtain the final bound
\[
\frac{1}{T}\sum_{t=0}^{T-1}\mathbb{E}\!\bigl[\|\nabla f(w_{t})\|^{2}\bigr]
\le
\frac{2\bigl(f(w_{0})-f^{\star}\bigr)}{L\,\bar\eta_{T}T}
+\frac{L\,\eta_{\max}\sigma^{2}}{2},
\]
which is exactly inequality (3).  ∎
\end{enumerate}

\section*{Rate Derivation (With Constants)}
From Theorem \ref{thm:main} we have
\[
\mathcal{G}_{T}:=
\frac{1}{T}\sum_{t=0}^{T-1}\mathbb{E}\!\bigl[\|\nabla f(w_{t})\|^{2}\bigr]
\le
\underbrace{\frac{2\Delta_{0}}{L\bar\eta_{T}T}}_{\text{deterministic term}}
\;+\;
\underbrace{\frac{L\eta_{\max}\sigma^{2}}{2}}_{\text{variance term}},
\qquad
\Delta_{0}:=f(w_{0})-f^{\star}.
\]

If we choose a cyclic schedule with a fixed period $M$ and constant
amplitudes $\eta_{\min},\eta_{\max}$, then
\[
\bar\eta_{T}= \eta_{\min}+\frac{\eta_{\max}-\eta_{\min}}{2}
\;+\;O\!\Bigl(\frac{1}{T}\Bigr),
\]
i.e. the average converges to a constant $\bar\eta:=\eta_{\min}+(\eta_{\max}-\eta_{\min})/2$.
Hence the deterministic term decays as $O(1/T)$.

To obtain the canonical $O(T^{-1/2})$ rate we set
\[
\eta_{\max}=c\,T^{-1/2},\qquad \eta_{\min}=c'\,T^{-1/2},
\]
so that $\bar\eta_{T}=O(T^{-1/2})$ and the variance term becomes $O(T^{-1/2})$ as well.
Consequently,
\[
\mathcal{G}_{T}=O\!\bigl(T^{-1/2}\bigr).
\]

\section*{Special Cases}
\begin{enumerate}[label=\alph*. , leftmargin=*]
\item \textbf{Constant Learning Rate.}  
If $\eta_{\min}=\eta_{\max}=\eta$, the cyclic schedule reduces to standard SGD
and (3) recovers the classical bound
$\frac{2\Delta_{0}}{L\eta T}+ \frac{L\eta\sigma^{2}}{2}$.

\item \textbf{Purely Increasing Schedule.}  
If we only use the first half of the triangle (i.e. $\eta_{t}$ linearly
increases from $\eta_{\min}$ to $\eta_{\max}$ and then stays at $\eta_{\max}$),
the same proof applies because the only property required is the uniform
upper bound $\eta_{\max}$.  The average $\bar\eta_{T}$ is larger,
hence the deterministic term improves, while the variance term remains
controlled by $\eta_{\max}$.

\item \textbf{Cosine Annealing.}  
A cosine schedule $\eta_{t}= \eta_{\min}+ \tfrac12(\eta_{\max}-\eta_{\min})
\bigl(1+\cos(\pi t/M)\bigr)$ also satisfies $ \eta_{\min}\le\eta_{t}\le\eta_{\max}$,
so the theorem holds verbatim.  The proof does not depend on the exact shape,
only on the uniform bounds.
\end{enumerate}

\end{document}